\documentclass[12pt]{article}
\usepackage[margin=1in]{geometry}
\usepackage{setspace}
\usepackage{hyperref}
\usepackage{amsmath,amssymb}

\hypersetup{
    colorlinks=true,
    linkcolor=blue,
    urlcolor=blue,
    citecolor=blue
}

\begin{document}
\singlespacing

\begin{flushright}
Prof. Humberto Bernal \\
Universidad Colegio Mayor de Cundinamarca \\
Ph.D. in Economics, Universidad de los Andes, Colombia \\
E-mail: \href{mailto:hbernal@uniandes.edu.co}{hbernal@uniandes.edu.co}, \href{mailto:hbernalc@universidadmayor.edu.co}{hbernalc@universidadmayor.edu.co} \\
ORCID: \href{https://orcid.org/0000-0002-4864-1726}{0000-0002-4864-1726} \\
\today
\end{flushright}

\vspace{1.5em}

\begin{flushleft}
The Editorial Board \\
\textit{Conflict Management and Peace Science}
\end{flushleft}

\vspace{1.5em}

\noindent Dear Editors of \textit{Conflict Management and Peace Science},

\vspace{1.5em}

\noindent It is a great privilege to submit the manuscript entitled \textbf{``Identifying Rational Types in Unknown Environments under an Indirect Dynamic Mechanism''} for consideration for publication in \textit{Conflict Management and Peace Science}.

This paper develops a novel dynamic mechanism design framework tailored to environments where information asymmetry, strategic secrecy, and operational conflict dynamics are paramount. Specifically, the framework addresses the critical challenge of screening and type identification of armed and criminal groups under incomplete information, using the empirical context of kidnapping for ransom during the armed conflict in Colombia.

The manuscript addresses a central issue in the formal study of conflict, deterrence, and non-state armed actors: how the state can design dynamic intervention policies (balancing financial interdiction $\alpha$ and operational pressure $\gamma$) when distinct perpetrator groups (guerrillas such as FARC and ELN, paramilitary organizations, and common criminal gangs) present observational pooling in their short-run demands. By formalizing the strategic interaction among the State, the Captor, and the Victim's Family through an indirect dynamic mechanism with endogenous execution noise (the Mano de Dios--Guadalupe process), the paper provides three key contributions:

\begin{enumerate}
    \item \textbf{A Dynamic Game of Conflict and Deterrence:} We formalize the interaction between the state, armed captors, and civil victims as an indirect dynamic mechanism, preserving the empirical realism of observable enforcement instruments rather than reducing them to abstract direct revelations.
    
    \item \textbf{Identification of Armed Perpetrator Types:} We establish the Rational Type Identification Theorem, proving that recurrently separated observable signals (captivity duration and competing-risk outcomes: rescue, ransom payment, escape, or execution) enable the state's Bayesian beliefs to converge almost surely to the true criminal technology of the captor in prolonged captivity.
    
    \item \textbf{Empirical Calibration and Policy Trade-offs:} Using comprehensive microdata from the Centro Nacional de Memoria Histórica (CNMH) in Colombia, we calibrate the model via competing-risks duration models. The structural simulations demonstrate that increasing operational pressure shortens captivity duration for insurgent groups by precipitating rescues, whereas for common crime it can inadvertently increase duration due to logistical dispersion constraints.
\end{enumerate}

To ensure complete transparency and reproducibility, all mathematical proofs are verified in the \textbf{Lean 4 Interactive Theorem Prover}, econometric scripts in Stata are documented, an interactive simulation portal is provided at \url{https://7kkgn4sfscykkfvv5794qw.streamlit.app/}, and all replication materials are available at \url{https://github.com/Donzhumber/CMP.git}.

Given the distinguished tradition of \textit{Conflict Management and Peace Science} in publishing cutting-edge formal models of conflict, deterrence, and quantitative peace research, we believe this manuscript will be of high interest to your readership.

Thank you very much for your time and consideration.

\vspace{2.5em}

\begin{flushleft}
Sincerely,

\vspace{2.5em}

\textbf{Prof. Humberto Bernal, Ph.D.} \\
Universidad Colegio Mayor de Cundinamarca \\
Doctor in Economics, Universidad de los Andes, Colombia
\end{flushleft}

\end{document}
